% !TeX spellcheck = de_DE

Die Evaluation des Proof of Concept erfolgt anhand verschiedener Textbeispiele aus dem Anwendungsbereich der LS+-Pipeline.
Es gibt dabei schon interne Texte, welche für Tests der bestehenden Pipeline verwendet wurden, bei diesen Texten gibt es auch generierte LS+ Texte vor dem Feedbackloop.
Um die Wirksamkeit des Ansatzes zu erörtern, werden mehrere Vergleichsszenarien betrachtet.

Zunächst werden die ursprünglichen Texte mit den durch den Feedbackloop überarbeiteten Versionen verglichen.
Darüber hinaus erfolgt ein Vergleich der bereits übersetzten LS+ Texten mit den überarbeiteten Versionen.
Dies wird den Hauptteil der Evaluation darstellen, da der Feedbackloop für diese Stelle angesetzt ist.
Zusätzlich werden manuell erstellte LS+-Texte als Referenz herangezogen und verglichen.

Die Bewertung erfolgt sowohl anhand der verwendeten Qualitätsmetriken als auch durch eine generelle Betrachtung der Kohärenz.
Dadurch soll untersucht werden, ob sich die durch den Feedbackloop erhoffte Verbesserungen nachweisen lassen und welche Grenzen der Ansatz aufweist.